\documentclass{article}
\usepackage[utf8]{inputenc}
\usepackage[T1]{fontenc}
\usepackage{geometry}
\geometry{
    a4paper,
    left=2.5cm,
    right=2.5cm,
    top=2.5cm,
    bottom=2.5cm
}
\usepackage{graphicx}
\usepackage{float}
\usepackage{hyperref}
\usepackage[english]{babel}
\usepackage{amsmath}
\usepackage{amssymb}
\usepackage{array}
\usepackage{booktabs}
\usepackage{caption}
\usepackage{subcaption}
\usepackage{enumitem}
\usepackage{fancyhdr}
\usepackage{titlesec}
\usepackage{url}
\usepackage{textcomp}
\usepackage{xcolor}
\usepackage{listings}
\usepackage{multirow}
\usepackage{longtable}
\usepackage{tabularx}
\usepackage{times}

% Set up page style
\pagestyle{fancy}
\fancyhf{}
\rfoot{\thepage}
\renewcommand{\headrulewidth}{0pt}

% Title formatting
\titleformat{\chapter}[display]{\normalfont\bfseries}{\chaptertitle\thechapter}{0pt}{\Huge}
\titlespacing*{\chapter}{0pt}{50pt}{40pt}

% Document info
\title{Template Document}
\author{Academic Support Agent}
\date{\today}

\begin{document}
\maketitle

%Define el fondo de la primer diapositiva
{%Inicia el cambio de fondo para la portada
\setbeamertemplate{background}{%
\includegraphics[width=\the\paperwidth]{images/FONDO_AL_15.png}
}
%lo anterior es para definir el fondo de la primer diapositiva

\begin{frame}
  \titlepage %Necesario para generar la portada
\end{frame}
}%aqu  termina el cambio de fondo

\begin{frame}

\end{frame}

\section{Introducci n} %%T tulo de la secci n (Opcional)
\subsection{Ejemplo} %%T tulo de la subsecci n

\begin{frame}
  \frametitle{T tulo de la diapositiva} 
  \framesubtitle{Subt tulo} %%Subt tulo de la diapositiva (opcional)
  Algo de texto en la parte superior
  
  \[
   B * \pi =\int_{0}^{\infty} \frac{5x^2}{\sqrt{a+b}}\, dx
   \]

  \begin{itemize}
    \item[\checkmark] Un elemento en la lista %%[\checkmark] muestra una palomita al inicio de la l nea
    \item Segundo elemento
    \item Otro elemento
    \item Y finalmente, otro m s
  \end{itemize}
\end{frame}

\section{Secci n de ejemplo 1}
\begin{frame}{Otra diapositiva} %%Otra forma (m s corta) de poner el t tulo a la diapositiva
Texto adicional. No tiene un prop sito en particular m s que ocupar algo de espacio.
Texto adicional. No tiene un prop sito en particular m s que ocupar algo de espacio.
Texto adicional. No tiene un prop sito en particular m s que ocupar algo de espacio.
Texto adicional. No tiene un prop sito en particular m s que ocupar algo de espacio.
\( \sum_0^{\infty} a_i \) 
Texto adicional. No tiene un prop sito en particular m s que ocupar algo de espacio.
\end{frame}

\section{Secci n de ejemplo 2} %%Otra secci n
\begin{frame}
    \frametitle{Un teorema}
    
    \begin{teorema} %%Uso del "environment" definido al inicio del documento.
        Los n meros primos son infinitos
    \end{teorema}

    Y algo de texto
    
    \centering \includegraphics[width=0.45 \paperwidth]{images/Manifold1.png}
\end{frame}

\section{Teorema con referencia}
\begin{frame}{Ejemplo}
\framesubtitle{Incluye referencia}
    \begin{teorema}[Teorema de Stokes]
        Sea $M$ una variedad compacta orientada con frontera, $\omega$ una $k-1$-forma sobre $M$, entonces $$\int_{M} d\omega = \int_{\partial M} \omega$$
    \end{teorema}

    la prueba se puede revisar en: \cite{spivak2018calculus}
\end{frame}

\section{Referencias}
\begin{frame}{Referencias}
    %posibles estilos para bibliograf a: unsrt , siam , plain , ieeetr , alpha , acm , abbrv
    
  
\end{frame}


\begin{frame}{Gracias}
    
\end{frame}
\end{document}